\section{Trajectory Data Blocks}\label{sec:trajectory-data-blocks}

Trajectory data blocks contain information associated with an entire trajectory.
Segment data blocks, discussed in \cref{sec:segment-data-blocks}, are one kind of
trajectory data block. The remaining trajectory-level blocks are listed in
\cref{tab:trajectory-data-blocks}.

\begin{table}[htbp]
  \centering
  \caption{Trajectory data blocks.}
  \label{tab:trajectory-data-blocks}
  \begin{tabularx}{0.82\linewidth}{@{}>{\ttfamily}lX@{}}
    \toprule
    \normalfont Data block & Description \\
    \midrule
    *define   & Defines user-defined output variables \\
    *dwn/crs  & Defines a downrange/crossrange reference \\
    *file     & Creates a column-oriented output file \\
    *iip      & Defines initial-impact-point variables \\
    *initial  & Defines initial conditions \\
    *print    & Selects printout variables \\
    *segment  & Defines a segment (\cref{sec:segment-data-blocks}) \\
    *tangent  & Defines the tangent-plane coordinate system \\
    \bottomrule
  \end{tabularx}
\end{table}

A trajectory begins with the \problemblock{*trajectory} keyword, a trajectory
number, a trajectory name, the keywords \texttt{start on}, and a segment number.
For example,

\begin{lstlisting}[style=taosmanual]
*trajectory 2 target start on 31
\end{lstlisting}

begins trajectory 2 for a vehicle called \texttt{target} on segment 31. Every item
following \problemblock{*trajectory} is required.

Trajectory information continues until another \problemblock{*trajectory} keyword
or a keyword unique to the problem level is encountered. Some keywords, including
\problemblock{*file} and \problemblock{*print}, are valid both inside and outside a
trajectory definition, so they do not terminate the trajectory. Only keywords that
are unique to the problem level, such as \problemblock{*trajectory},
\problemblock{*search}, and \problemblock{*units/fmt}, end a trajectory definition
(\cref{sec:problem-data-blocks}).

Every trajectory definition must contain an \problemblock{*initial} block giving the
initial conditions and at least one \problemblock{*segment} block. It normally also
contains at least one \problemblock{*print} block so that trajectory results are
written to the printout file. The following outline illustrates the organization:

\begin{lstlisting}[style=taosmanual]
*trajectory 1 missile start on 1
  *initial ...
  *print ...

  *segment 1 First segment definition
    ...
  *segment 2 Second segment definition
    ...

*trajectory 2 target ...
\end{lstlisting}

Indentation shows the hierarchy. Trajectory 1 extends from the first
\problemblock{*trajectory} keyword to the second. Between those keywords, initial
conditions are given in the \problemblock{*initial} block, printout variables are
listed in the \problemblock{*print} block, and segments are defined in the
\problemblock{*segment} blocks. Each \problemblock{*segment} block contains the
additional blocks described in \cref{sec:segment-data-blocks}, which tell TAOS how
to calculate the segment.

The following subsections give the details of each trajectory data block.
